\section{The Trivium \& Mystery}


Cologero's translations\footnote{\url{https://gornahoor.net/?p=6649}} have provided this gem, from de Giorgio:

\begin{quotex}
We could also call it ``intuition'' although no psychological quality is given to this term: the psyche in fact is below the spirit, the intellect, the heart---these three terms denoting, under three aspects, the same type of integrative activity of the divine. The spirit expresses the direct integration whose absolute type is the divine breath, the intellect expresses the cognitive permeation, the heart expresses radiant receptivity: by means of the first, one is elevated, with the second, absorbed, in the third, one is welcomed and realizes himself. Representing here a vertical axis, the spirit is the peak, the intellect the base, the heart the center that gathers the two extreme points and extends them, prolonging them horizontally, hence the Cross as radiant symbol of universality and unifying centrality.
\end{quotex}


The intellect, the heart, and the spirit are One. I would like to relate this to the Trivium\footnote{\url{https://en.wikipedia.org/wiki/Trivium}}, in the Christian Western Tradition. The Trivium consists of grammar (which is emphasized first), then logic (which crops up second, as the pupil reaches puberty, which is age-appropriate), and finally, rhetoric, where the student begins to be the master of his knowledge, rather than being eternally dominated by Facts in a kind of Gradgrindism\footnote{\url{https://en.wikipedia.org/wiki/Gradgrind}}. Now, none of these occur without the other -- even young children learn a form of rhetoric, in the form of songs and poetry and stories. It is merely that the emphasis of the logical progression occurs in an age more suited for it, while the other two categories receive only their due. If we were relating this to a knowledge of castes and Order (including internal order) we would say that the Spirit correlates to Fact, the Heart to Logic, and the Intellect to Rhetoric (for true self-consciousness is consciousness of the Master inside, allowing domination of both Self and Destiny). The Spirit initiates, or factualizes the supernatural; the Heart has its inner logic that we must learn to listen to; the Intellect realizes and epitomizes the power of the Invisible made Visible, the fruit of what has gone before.

This explains, by the way, the strange association of heart \& head, the primacy of God's calling, and the divinization of man within Supernature, becoming God's equal by Grace.

I did not read this schema -- it was obvious to me, latently so, when I read the passage above from de Giorgio. On the contrary, one might choose to say, at times, that the Spirit leads by an inner logic, the Heart creates a vision or rhetoric of the goal, and the Intellect deals with the ``facts'' the matter. I would reply -- the student/pilgrim will at various stages emphasize different portions of their Trivium, and so a change of emphasis or the seeming appearance of a ``shuffle'' of responsibility in certain circumstances, does not destroy the analogy. The exception, in other words, should strengthen the rule: the analogy fits.

The Trivium was made explicit during the Carolingian period\footnote{\url{https://en.wikipedia.org/wiki/Carolingian\_era}} of Western civilization, and is surely cognate in some inner way to the doctrine of Trinity; if nothing else, reflection such as Augustine's on the tripartite nature of the soul and its relation to the tripartite functions of the persons of God encouraged a numerical attention to Three. Adding this to the Quadrivium produces a perfect number, Seven\footnote{\url{https://gornahoor.net/?p=5879}}. There were, of course, ancient Pythagorean doctrines (derived from Egypt) which taught ruler-ship by the Hebdomads or Sevens. It should be unnecessary to point out that the medievals would have seen significance in Four Gospels, Three Persons of the Trinity, Two Natures in One Christ, and the Divine Unity.

Are the ``Three'' related to\footnote{\url{http://biblehub.com/1\_john/5-8.htm}} the Spirit, the Water, and the Blood of Saint John?

\begin{quotex}
And there are three that bear witness in earth, the spirit, and the water, and the blood: and these three agree in one.
\end{quotex}


One of the discouraging things I see in Christianity is an inability or unwillingness to reconsider previous meditations in light of new ones. Are the liberals the only ones who are allowed to be daringly adventurous, since the Fathers who wrote the \emph{Philokalia} have not their like today?

For instance, we might ask ourselves, as Christians, more than just where the revelations to Adam or Enoch went. We could even ask more than whether Guenon was right to argue that Christianity was just one vehicle or aspect of the Tradition, or whether men like Jean Borella\footnote{\url{http://www.religioperennis.org/documents/Editorial/Issue5/editorialIII1.pdf}} are right to criticize Guenon as a ``Judaizer'' of sorts (the answer to this is mostly irrelevant to the spiritual level that most Christians, or most people, for that matter, would find themselves at in this life). My belief is that, in looking for correspondences in other religions (if done in obedience to the Spirit), Christians will be strengthened and surprised in ways that intensify the experience of their own Tradition, if it is kept rich enough. Who, for instance, is the Sophia\footnote{\url{http://en.wikipedia.org/wiki/Sergei\_Bulgakov}} in Proverbs? We have an advantage in our Era, in that, certain questions having been tentatively elucidated and put forward, the modern situation is such that no conclusively binding regulation on conscience has been reached by the Faith. Perhaps we are missing great opportunities here?

I will suggest something else for those who are interested in esoteric Christianity. During Pentecost, each person heard the sermon of Peter \emph{in their own tongue}. Since Pentecost only recapitulates what had already happened I would like to posit, as a means of reconciling the findings of Nag Hammadi with Orthodoxy, that Jesus' inner circle of \emph{discipulos} each heard the Christ speaking in the language of their soul certain secret sayings known only to them. Whether this occured at once, or whether it was done mystically, or one at a time, we take it at face value what the Gospels and other texts confess: that these are the sayings of Jesus \emph{given to this or that one} of the disciples\footnote{\url{http://gnosis.org/naghamm/nhl\_thomas.htm}}. In other words, each was given the ``One'' in a form that they could understand and pass on. The Synoptics look very much the same, book to book, but try reading Mark and then skipping to John.

This presents no problem for Christian theology, in that there may still be preserved four exoteric Gospels, perhaps (also) for that reason kept more accurately and therefore more reliably. This presents no problem for esotericism within Christianity, as it is not dependent on the Letter or written record in any case. But it does reconcile what we know about history and archaeology, in broad enough terms, to declare a peace with Tradition.

If Christians were more attentive to our own Trivium, we would have no need of Traditionalists to teach it to us, today. Or Gnostics, for that matter, real or imagined.


\flrightit{Posted on 2013-07-19 by Logres}

